% =====================================================================
%  Chapter 9 — Principle: Accountability (Make Danger Visible)
%  Map: research/chapter-map.md (Chapter 9)
%  Shape: idea → origin → cost → what it opens up. Ideas, not how-to.
%  Rule: the author is never quoted in body text; short quotes only in
%  footnotes. Full statements live in interviews/.
% =====================================================================
\chaptersubtitle{Make Danger Visible}
\chapter{Accountability}
\label{ch:danger}

\epigraph{An erroneously implemented native function can cause a VM internal
state inconsistency, which can cause a crash of the VM, or miscellaneous
misbehaviors of the VM at any point after the call to the native function.}
{Erlang/OTP documentation, \code{erl_nif}}

No useful language can refuse all contact with the outside world. Operating
systems, device drivers and decades of trusted libraries are written in C. The
question is not whether a program will ever touch unsafe code. It is whether
anyone responsible for the program will know.

\section{Unsafe impedance}

The 2024 paper that opens this book's story (Chapter~\ref{ch:why}) measured
languages by how hard they make it to write, load and hide unsafe code. Its
complaints were mostly about visibility. Some languages gave no sign at all
that a piece of code relied on something unsafe. Others gave a sign too quiet
for a newcomer to notice. Rust's own guidance encouraged libraries to wrap
unsafe internals in a safe-looking surface. Even Erlang, the paper's
best-rated platform, let ordinary code call native functions with no warning
that it was doing so.

The paper's remedy was a human process: no engineer could add native code
alone, and a deliberately skeptical review had to approve it. It also admitted
the limit of any such process. People can make unsafe code rarer; they cannot
make it impossible.\footnote{``Business practices such as UAPs can not
guarantee code safety, only reduce the probability.''}

\section{From process to compiler}

Silica moves that review into the language. Its bridge to foreign code is
named Fifi, after a poodle that looks cuddly and bites when you pet
it.\footnote{\src{silica/docs/index.md}, line 40.} Any part of a program that
uses foreign code must say so in its very name, with the prefix
\code{dangerous_}, and so must every part that uses \emph{that}, all the way up to
the program as a whole. Danger cannot be wrapped, inherited quietly or buried in
a dependency. Putting the marker in names, rather than among the other effects a
piece of code declares, keeps it pure signal that can never be lost in noise. Foreign code runs only in separate workers set aside for it, data
it returns is treated with suspicion, and other languages cannot call into
Silica at all. If the foreign code crashes, only its worker fails, and its
supervisor restarts it (Chapter~\ref{ch:failsoft}).

Fifi does not claim to make foreign code safe. Its design says plainly that it
is not a sandbox. What it promises is that the boundary cannot be hidden.

\section{A bridge, reluctantly built}

Fifi was not part of the original plan. Silica's founding notes refused any
connection to C at all, and Barney did not want to add one. What changed the
decision was a picture of the people Silica most needed to reach: the academics
and professionals who write the code that runs directly on boards and bare
computers. Those boards are everywhere. They control cars, electrical and water
systems, networks, and most of the computers people touch every day.

Barney knew many of these programmers, and knew how deep their roots in C run,
along with the habit of picturing hardware as the PDP-11 once was. Persuading
them even to look at another language would be long, hard work, and so far, for
most of them, it has been. A language that demanded they abandon every line of C
at once would never be tried. Fifi lets them move over gradually, keeping the
foreign code they cannot yet replace, while making its danger impossible to
miss.\footnote{Interview, 2026-11-13.}

The same honesty that shaped the 2024 paper shaped this decision. Barney could
not claim that the BEAM languages were perfect, and does not claim it of Silica
either. Fifi is a compromise, and Silica says so in the name of every module that
uses it.

\section{What SEMP left behind}

The BEAM\_SEMP project of 2025 spent two months removing dangerous defaults
from a running Erlang system: stopping peers from creating new names, from
loading code, from calling functions remotely. In Silica, most of those
defaults never exist in the first place. There was nothing to remove.

% TODO-INTERVIEW: SEMP judged peers by their behaviour over time. Silica has
%   no counterpart yet. Dropped, or waiting?

\section{Authority you can see}

The direction of travel goes further: a program should be able to do only what
the parts it has been given are able to do, and that authority should be
visible in the design. Much of this is planned rather than built, and the
documents say so.

% TODO-VERIFY: Keep status current: capabilities, brokered isolation, TLS and
%   the crypto proposal were design-only as of 2026-11-13.

\section{What it costs}

Visibility is uncomfortable. Using one foreign library renames everything
above it, and that is the point: the inconvenience is the signal. Some things
that are easy in other languages, such as letting outside code call in, are
simply not offered.

\section{What it opens up}

\begin{itemize}
  \item \textbf{Audits by search.} Whether a release contains foreign code is
        answered by searching for one word, not by reverse-engineering its
        build.
  \item \textbf{Supply chains that show themselves.} A dependency that brings
        in unsafe code cannot arrive unannounced; its presence changes the
        names of everything that relies on it.
  \item \textbf{Policies that can be enforced.} An organisation can decide
        that a product contains no foreign code, or only approved foreign
        code, and check the decision mechanically.
  \item \textbf{An answer to the regulators' question.} Governments now ask
        software makers to be secure by design. A language that makes danger
        visible gives them evidence rather than assurances.
\end{itemize}

% TODO-INTERVIEW: Who do you most want to read this chapter — engineers,
%   managers, or the people writing security policy?
